\Needspace{8\baselineskip}\originalchapter{2002 年 Landau 响应与碰撞算子作业}
\originalsection{2002 PS7：屏蔽与两类速度空间不稳定}
\begin{exercise}\label{ex:hist02-7-1}
PS7 原题1（有官方答案）。在初始热平衡、无碰撞的等离子体于 $t=0$ 突加静止点测试电荷 $q_T$，忽略离子运动并取线性响应。用题给电子介电函数 $\epsilon(k,\omega)=1+[1+\zeta Z(\zeta)]/(k^2\lambda_{De}^2)$，$\zeta=\omega/(kv_{te})$，写 Fourier--Laplace 势，取 $t\to\infty$ 并与电子 Boltzmann 响应比较；说明初值项何时消失。
\end{exercise}
\begin{solution}
编者独立解答。以 $\nabla^2\phi=-\rho/\epsilon_0$ 和课程 $e^{\ii\omega t}$ Laplace 约定，测试源的时间变换为 $\ii q_T/\omega$，故
\begin{equation}\phi_{k\omega}=\frac{\ii q_T}{\epsilon_0\omega k^2\epsilon(k,\omega)}+\frac{S_0(k,\omega)}{k^2\epsilon(k,\omega)}.
\end{equation}
$S_0$ 由插入瞬间已有粒子扰动决定；若初值不持续驱动且 Landau 稳定，逆变换中它的阻尼极点与相混合贡献渐消。终值定理取 $\epsilon(k,0)=1+1/(k^2\lambda_{De}^2)$，有 $\phi_{k,\infty}=q_T/[\epsilon_0(k^2+\lambda_{De}^{-2})]$，逆 Fourier 得 $\phi_\infty(r)=q_Te^{-r/\lambda_{De}}/(4\pi\epsilon_0r)$。Boltzmann 线性化 $\delta n_e=n_0e\phi/T_e$ 代 Poisson 式给同一 Helmholtz 核；这说明静态屏蔽无需把碰撞造成的真实热平衡误认为必须条件。若有不衰减模、驱动初值或不稳定分布，不能直接丢 $S_0$。
\end{solution}

\begin{exercise}\label{ex:hist02-7-2}
PS7 原题2（有官方答案）。原件给沿波矢方向的速度分布（各向同性横向速度已积出）：\begin{equation}\bar f_e(u)=\frac{n_p+n_b}{\sqrt\pi v_e}e^{-u^2/v_e^2}+\frac{n_b}{\sqrt\pi v_b}e^{-(u-u_b)^2/v_b^2}.\end{equation}原件的三维式背景系数确为 $n_p+n_b$；在 $n_b\ll n_p$ 首阶可用 $n_p$。另给 $u_b\gg v_e$、$v_b\sim v_e$，离子不动。画分布并指出 bump-on-tail 增长位于哪侧、频率与波数范围，求最快模的 $\omega_r,k,\gamma$ 领先阶。
\end{exercise}
\begin{solution}
编者独立解答。背景电子使高频 Langmuir 实根 $\omega_r\simeq\omega_{pp}$；束尾左坡 $u\lesssim u_b$ 有 $\partial f_b/\partial u>0$，故逆 Landau 阻尼可胜过背景在 $u\gg v_e$ 的指数小阻尼。设 $\xi=(\omega_r/k-u_b)/v_b$；领先束密度阶的介电虚部给
\begin{equation}\gamma\simeq-\sqrt\pi\frac{\omega_{pb}^2\omega_{pp}}{k^2v_b^2}\xi e^{-\xi^2},\quad \omega_{pb}^2/\omega_{pp}^2=n_b/n_p.
\end{equation}
正增长要求 $\xi<0$，主要在 $k\sim\omega_{pp}/u_b$ 附近，带宽 $\Delta k/k\sim v_b/u_b$；固定慢变的 $k^{-2}$ 因子时 $-\xi e^{-\xi^2}$ 在 $\xi=-1/\sqrt2$ 最大。因此 $\omega_r\simeq\omega_{pp}$、$k_{\max}\simeq\omega_{pp}/(u_b-v_b/\sqrt2)$，$\gamma_{\max}\simeq\sqrt{\pi/(2e)}(n_b/n_p)(u_b/v_b)^2\omega_{pp}$。该微扰答案还需 $\gamma\ll\omega_{pp}$ 和束背景阻尼可忽略；若 $(n_b/n_p)(u_b/v_b)^2$ 不小，须解完整色散式，不能套此最速值。
\end{solution}

\begin{exercise}\label{ex:hist02-7-3}
PS7 原题3（有官方答案）。准中性电子—质子等离子体：离子静止 Maxwell、电子以 $u_d$ 漂移 Maxwell，$v_i<u_d<v_e$、$v_s=\sqrt{2T_s/m_s}$，三维分布分别正比 $e^{-u^2/v_i^2}$ 和 $e^{-|\mathbf u-u_d\hat z|^2/v_e^2}$，总密度均为 $n$。画同轴速度剖面；在 $v_i<\omega_r/k\ll v_e$ 的离子声范围写含电子、离子贡献的 $\gamma/\omega_r$，指出 $T_e\gg T_i$ 的作用和可能的频率／波数区间。
\end{exercise}
\begin{solution}
编者独立解答。$k\lambda_{De}\ll1$ 时 $\omega_r\simeq kc_s$，$c_s=\sqrt{T_e/m_i}$；令 $\xi_e=(c_s-u_d)/v_e$、$\xi_i=c_s/v_i=\sqrt{T_e/(2T_i)}$。用 $\gamma=-\operatorname{Im}\epsilon/(\partial_\omega\operatorname{Re}\epsilon)$，在 $|\xi_e|\ll1$、$\xi_i\gg1$ 的首阶得到
\begin{equation}\frac\gamma{\omega_r}\simeq-\frac{\sqrt\pi}{2}\left[\frac{c_s-u_d}{v_e}+\frac{T_e}{T_i}\xi_i e^{-\xi_i^2}\right].\end{equation}
第一项在 $u_d>c_s$ 时为驱动，第二项是总为正的离子 Landau 阻尼，故还需 $(u_d-c_s)/v_e>(T_e/T_i)\xi_i e^{-\xi_i^2}$。$T_e\gg T_i$ 压小离子尾的指数因子，而 $T_e\sim T_i$ 时离子阻尼强、上述大 $\xi_i$ 渐近式也失效。可能范围是 $\omega_r\ll\omega_{pe}$、$0<k\lambda_{De}\ll1$ 且相速落在 $v_i<c_s<u_d<v_e$；给定题面没有单独的 $u_d,T_e,T_i$ 数字，不应虚构唯一 $k$ 数值。
\end{solution}

\originalsection{2002 PS8：Landau 算子守恒和温度平衡}
\begin{exercise}\label{ex:hist02-8-1}
PS8 原题1（有官方答案）。同种粒子的 Landau 算子 $C(f,f)=\Gamma\nabla_v\cdot\int U(\mathbf v-\mathbf v')[f'\nabla_vf-f\nabla_{v'}f']\dd^3v'$，$U(\mathbf u)=(I-\mathbf u\mathbf u/u^2)/u$，题卷用 Gaussian 单位 $\Gamma=2\pi q^4\ln\Lambda/m^2$。证明粒子数、总动量和总动能各自不变。
\end{exercise}
\begin{solution}
编者独立证明。数目积分是速度散度的边界项，若分布衰减足够快则为零。记对称双积分核 $\mathbf A=U(\mathbf u)(f'\nabla_vf-f\nabla_{v'}f')$，交换 $\mathbf v\leftrightarrow\mathbf v'$ 得 $\mathbf A\mapsto-\mathbf A$。动量变化 $-m\Gamma\iint\mathbf A\dd^3v\dd^3v'=0$。能量变化 $-m\Gamma\iint\mathbf v\cdot\mathbf A\dd^3v\dd^3v'=-m\Gamma\iint(\mathbf v-\mathbf v')\cdot\mathbf A\dd^3v\dd^3v'/2=0$，因为 $\mathbf u\cdot U(\mathbf u)=0$。所以每次同种二体散射只在速度空间重排，而不改这三个碰撞不变量。
\end{solution}

\begin{exercise}\label{ex:hist02-8-2}
PS8 原题2（有官方答案）。交叉 Landau 算子的核仍为上题 $U$，但 $C_{\alpha\beta}=\Gamma_{\alpha\beta}\nabla_v\cdot\int U(\mathbf v-\mathbf v')[(m_\beta/m_\alpha)f_\beta'\nabla_v f_\alpha-f_\alpha\nabla_{v'}f_\beta']\dd^3v'$，原件 Gaussian 系数 $\Gamma_{\alpha\beta}=2\pi Z_\alpha^2Z_\beta^2e^4\ln\Lambda/(m_\alpha m_\beta)$。证明每个物种分别守恒粒子数，但动量与能量只在 $C_{\alpha\beta}+C_{\beta\alpha}$ 的物种和中守恒；解释个别物种可交换什么。
\end{exercise}
\begin{solution}
编者独立证明。两个 $C$ 均为各自速度散度，故 $\int C_{\alpha\beta}\dd^3v=\int C_{\beta\alpha}\dd^3v'=0$。交叉核满足作用反作用的对称化：把 $\alpha$ 的 $m_\alpha\mathbf v$ 矩与 $\beta$ 的 $m_\beta\mathbf v'$ 矩相加，交换物种与两速度后积分反号相消；能量矩相加后剩余项正比 $(\mathbf v-\mathbf v')\cdot U(\mathbf v-\mathbf v')=0$。单独如 $\int m_\alpha\mathbf v C_{\alpha\beta}$ 或 $\int m_\alpha v^2C_{\alpha\beta}/2$ 一般不为零：碰撞使两物种交换漂移动量与热能，直到共同漂移、温度。这个“和”是两种算子合计，不能误解为每种各自守恒。
\end{solution}

\begin{exercise}\label{ex:hist02-8-3}
PS8 原题3（有官方答案）。电子、离子初为各自 Maxwell，题给 $C_{ei}\simeq C^L+C^E$，其中 $C^L=\nu_{ei}(v_e/v)^3\frac12\partial_\mu[(1-\mu^2)\partial_\mu f]$，$C^E=\nu_{ei}m_ev_e^3/(2m_i)\{v^{-2}\partial_v+[T_i/(m_ev^2)]\partial_v(v^{-1}\partial_v)\}f$。取 $m_ev^2/2$ 矩，推出题卷目标式 $\partial_t(3nT_e/2)=-\nu_{ei}(m_e/m_i)\sqrt{2/\pi}\,n(T_e-T_i)$，比较散射角与热平衡速率。此题明写 $v_e=\sqrt{T_e/m_e}$，与 PS7 的 $\sqrt{2T_e/m_e}$ 不同；还须核对目标式和印刷算子的自洽性。
\end{exercise}
\begin{solution}
编者独立解答。$C^L=\nu_{ei}(v_e/v)^3\frac12\partial_\mu[(1-\mu^2)\partial_\mu f]$，对各向同性 Maxwell 的 $\partial_\mu f=0$，并且即便各向异性，角积分也没有能量矩。能量交换只来自题给径向算子\begin{equation}C^E(f)=\nu_{ei}\frac{m_e v_e^3}{2m_i v^2}\frac{\partial}{\partial v}\left[f+\frac{T_i}{m_ev}\frac{\partial f}{\partial v}\right].\end{equation}对规范化 Maxwell 分布，$\partial_vf_e=-(m_ev/T_e)f_e$，方括号恰是$(1-T_i/T_e)f_e$，而不是多一个$1/v$。取能量矩并对$v$分部积分后，$\int_0^\infty vf_e\dd v=n/[(2\pi)^{3/2}v_e]$，得到
\begin{equation}\frac32n\frac{\dd T_e}{\dd t}=-\frac1{\sqrt{2\pi}}\frac{m_e}{m_i}\nu_{ei}n(T_e-T_i).
\end{equation}
原题目标式印为 $C_T=\sqrt{2/\pi}$。但若逐字采用同页写出的 $C^E$、$v_e^2=T_e/m_e$ 与规范化 Maxwell 分布\begin{equation}f_e(v)=\frac{n}{(2\pi v_e^2)^{3/2}}\exp\!\left(-\frac{v^2}{2v_e^2}\right),\end{equation}径向分部积分给 $C_T=1/\sqrt{2\pi}$；印刷目标差一个因子 2。扫描官方解答又把 Maxwell 指数写成 $e^{-v^2/v_e^2}$，却同时用了 $(2\pi v_e^2)^{-3/2}$ 的归一化，末行还将应有的 $T_i-T_e$ 误写成加号。因此本稿保留题面目标与逐字算子结果的区别，不冒称这些扫描式彼此严格一致。无论规范，角散射率 $\sim\nu_{ei}$，温度平衡率 $\sim(m_e/m_i)\nu_{ei}$；电子改方向远快于交换热能。
\end{solution}
